\documentclass[11pt]{article}
\usepackage[T1]{fontenc}
\usepackage[utf8]{inputenc}
\usepackage{lmodern}
\usepackage[margin=1in]{geometry}
\usepackage{amsmath,amssymb,amsthm,booktabs,longtable,array}
\setlength{\emergencystretch}{1em}
\usepackage{xcolor}
\usepackage{xurl}
\usepackage[colorlinks=true,linkcolor=black,citecolor=blue,urlcolor=blue]{hyperref}
\hypersetup{pdftitle={An Exact Local Extension and Homothetic Classes for OEIS A230543},
  pdfauthor={Carlo Corti}, pdfsubject={Arithmetic derivatives and homothetic classes},
  pdfkeywords={OEIS A230543, arithmetic derivative, exact computation},
  pdfcreator={LaTeX with hyperref}}
\newcommand{\D}{\operatorname{D}}
\newcommand{\Lrate}{\operatorname{L}}
\newcommand{\NN}{\mathbb{N}}
\newtheorem{theorem}{Theorem}[section]
\newtheorem{proposition}[theorem]{Proposition}
\newtheorem{corollary}[theorem]{Corollary}
\newtheorem{lemma}[theorem]{Lemma}
\theoremstyle{definition}
\newtheorem{definition}[theorem]{Definition}
\title{An Exact Local Extension and Homothetic Classes for OEIS A230543}
\author{Carlo Corti\\Independent researcher\\\texttt{carlo.corti@outlook.com}}
\date{October 2026}
\begin{document}
\maketitle

\begin{abstract}
OEIS A230543 lists the positive integers whose first two arithmetic
derivatives, together with the original integer, form the ordered legs of
an integer Pythagorean quadruple. Starting from an accepted 22-term
baseline, a complete C enumeration through $130000000000$ gives 18 further
terms and a 40-term local prefix. This indexing relies on the accepted
baseline and the program's negative coverage; a second independent
enumeration of all negatives was not performed. All 40 positive witnesses
were independently checked. We give a complete, elementary description of
the integer scales preserving any admissible primitive ordered quadruple.
The ten forms in the prefix comprise four infinite homothetic classes and
six singleton classes. Explicit proofs identify the four infinite families,
their narrower specializations and the singleton scales. Finite seed and
prime-pattern studies are reported separately from global classification.
The first unresolved term satisfies
$130000000000<a(41)\le244140625000$, with an exact upper witness but no
proved index for that witness. Historical logarithmic-derivative theory is
used with conservative attribution.
\end{abstract}

\section{The sequence and the scope of the extension}\label{sec:intro}
For nonnegative integers, the arithmetic derivative is defined by
$\D(0)=\D(1)=0$ and
\begin{equation}\label{eq:derivative}
 \D(n)=\sum_{p^e\parallel n}e\frac{n}{p}\qquad(n\ge2),
\end{equation}
where $p$ ranges over primes. Its product rule and logarithmic derivative
are classical~\cite{ufnarovski,kovic}. Throughout,
$\D^2(n)=\D(\D(n))$ denotes composition, not a square.
OEIS A230543~\cite{sequence} lists in increasing order, with offset one,
the integers $n\ge2$ for which $\D(n),\D^2(n)>0$ and
\begin{equation}\label{eq:predicate}
 n^2+\D(n)^2+\D^2(n)^2=R^2,\qquad R\in\NN_{>0}.
\end{equation}
We call $W(n)=(n,\D(n),\D^2(n),R)$ a witness. No primitivity,
coprimality, increasing order of the three legs, or membership in a
particular family is required.

The accepted OEIS baseline acquired on October 6, 2026 contains 22 terms,
ending at $a(22)=408918816$. The entry is due to Paolo P. Lava, with
terms 16--18 contributed by Giovanni Resta and terms 19--22 by Ray
Chandler~\cite{sequence}. The present study adds 18 terms relative to that
dated baseline. Attempts to retrieve the current A230543 internal entry
and b-file during the final checks were unsuccessful; we therefore do not
assert that all 18 are unpublished in the current public entry.

The computational statement is a conditional exact extension: the
accepted earlier indexing and complete negative coverage by the audited C
program imply the indices displayed in Table~\ref{tab:prefix}.
Independent trial-division checks establish membership of every positive
witness, but do not independently reproduce the full negative enumeration.
Section~\ref{sec:extension} separates these evidence scopes explicitly.

The structural contribution is a self-contained scale criterion for a
\emph{given} ordered primitive form, followed by its application to all
ten forms observed in the prefix. Historical integer-rate theory is
credited below. Ufnarovski and \AA hlander's Theorem 18~\cite{ufnarovski} already
solves rational logarithmic-derivative equations and gives the squarefree
denominator criterion. We restate its exponent decomposition in a
positive-integer form, with canonical residues and nonnegative exponent
increments, without a claim of absolute novelty. The
Hoepli and Sandit books of Balzarotti and Lava~\cite{book,sandit} were
not available in full for this study, limiting priority assessment.
The overview by Tossavainen et al.~\cite{overview}, now consulted in
full through an author-uploaded copy, surveys the history, elementary
bounds and differential equations of the arithmetic derivative.
No theorem is attributed to an unread book.

\section{Historical tools and ordered primitive forms}\label{sec:tools}
For $x>0$ put $\Lrate(x)=\D(x)/x$. Prime exponents in a product add,
so~\eqref{eq:derivative} gives, without a coprimality hypothesis,
\begin{equation}\label{eq:product}
 \D(xy)=x\D(y)+y\D(x),\qquad
 \Lrate(xy)=\Lrate(x)+\Lrate(y).
\end{equation}
The derivative is not additive: $\D(2)+\D(3)=2$, whereas $\D(5)=1$.
For a prime $p$, $\D(p)=1$ and $\D^2(p)=0$; moreover
$p^2<p^2+1<(p+1)^2$. Thus primes cannot satisfy the predicate.

Kovi\v{c}'s Proposition 2~\cite{kovic} describes the positive integer
solutions of $\D(t)=kt$, for an integer $k\ge1$, as
$t=\prod_{j=1}^k p_j^{p_j}$, with repetitions permitted. The case $k=0$
has only $t=1$. In particular, the fixed points $\D(t)=t$ are precisely
$p^p$, recorded in A051674~\cite{derivatives}. The 2024 survey
also documents the fixed-point question in the 1950 Putnam competition
and Barbeau's 1961 derivative bounds~\cite{overview}. These historical
scale factors will occur repeatedly.

For an integer $\lambda\ge0$ define
$T_\lambda(z)=\D(z)+\lambda z$ on nonnegative integers.
If $\Lrate(t)=\lambda$, the product rule gives
\begin{equation}\label{eq:iterate}
 \D^j(tm)=tT_\lambda^j(m)\qquad(j\ge0).
\end{equation}
Indeed the claim holds at $j=0$; if $z=T_\lambda^j(m)$ is the current
integer iterate, then $\D(tz)=t(\D(z)+\lambda z)$ establishes the next
step. For rational $\lambda$ the same induction is valid along any
chain whose displayed $T_\lambda$ iterates are integers; we do not apply
an integer-domain operator to an unspecified noninteger iterate.
This is product-rule induction, not a binomial expansion of an additive
operator.

\begin{definition}
For a witness $W=gQ$, let
$g=\gcd(n,\D(n),\D^2(n),R)$ and
$Q=(m,u,v,w)=W/g$. The ordered quadruple $Q$, with
$\gcd(m,u,v,w)=1$, is its primitive form. Its homothetic class is the set
of witnesses $tQ$ with $t$ a positive integer.
\end{definition}
The normalization uses all four coordinates and preserves their order.
The primitive form need not itself be a derivative witness.
By the product rule,
\begin{equation}\label{eq:shape}
 \lambda=\Lrate(g)=\frac{u-\D(m)}{m},\qquad
 \D(u)+\lambda u=v,\qquad m^2+u^2+v^2=w^2.
\end{equation}

\section{A complete scale criterion for a fixed form}\label{sec:scale}
\begin{theorem}\label{thm:shape}
Let $Q=(m,u,v,w)$ be the primitive form of a witness and let
$\lambda$ be given by~\eqref{eq:shape}. A positive integer $t$ produces
a witness $tQ$ if and only if $\Lrate(t)=\lambda$.
No noninteger positive rational multiplier produces four integer
coordinates from $Q$.
\end{theorem}
\begin{proof}
The equality $\D(tm)=tu$ is equivalent, by~\eqref{eq:product}, to
$\D(m)+m\Lrate(t)=u$ and hence to the stated rate condition.
If it holds, $\D(tu)=t(\D(u)+\lambda u)=tv$ and
$(tm)^2+(tu)^2+(tv)^2=(tw)^2$. All coordinates are positive. Also $tm\ge2$: otherwise $t=m=1$
and the first derivative equality would give $u=\D(1)=0$, a
contradiction. This proves necessity and sufficiency. If $a/b$ is reduced and
$(a/b)Q$ is integral, then $b$ divides every coordinate of $Q$, so $b=1$.
\end{proof}

The squarefree-denominator criterion and the decomposition into an
integer part plus prime-denominator fractions are classical
\cite[Theorem 18]{ufnarovski}. The following statement explicitly
restricts the exponent increments to those giving positive integer scales.

\begin{theorem}\label{thm:residues}
Suppose $\lambda=\Lrate(g)=A/B$ is reduced, with $B>0$. Then
$\lambda\ge0$ and $B$ is squarefree. For each prime $p\mid B$ let
$r_p\in\{1,\ldots,p-1\}$ be the representative satisfying
\begin{equation}\label{eq:residue}
 r_p\equiv A(B/p)^{-1}\pmod p.
\end{equation}
Put
\begin{equation}\label{eq:core}
 t_0=\prod_{p\mid B}p^{r_p},\qquad
 \alpha_0=\sum_{p\mid B}\frac{r_p}{p},\qquad
 k=\lambda-\alpha_0.
\end{equation}
For $B=1$ use $t_0=1$ and $\alpha_0=0$. Then $k$ is a nonnegative
integer, and all and only the positive integer solutions of
$\Lrate(t)=\lambda$ are
\begin{equation}\label{eq:allscales}
 t=t_0\prod_{j=1}^{k}p_j^{p_j},
\end{equation}
where the primes $p_j$ are arbitrary, may repeat, and may divide $B$.
\end{theorem}
\begin{proof}
Writing $g=\prod p^{e_p}$ gives $\Lrate(g)=\sum e_p/p\ge0$.
This sum has a denominator dividing the squarefree product of its
distinct primes, so its reduced denominator is squarefree.
For any proposed $t=\prod p^{f_p}$, let $M$ be the squarefree product
of the primes dividing $tB$. The rate equation is equivalent to
\begin{equation}\label{eq:modular}
 \sum_{p\mid t}f_p M/p=A M/B.
\end{equation}
If $p\mid B$ were absent from $t$, the left side would be zero modulo
$p$, whereas the right side would be nonzero, since $\gcd(A,B)=1$.
Thus every prime of $B$ divides $t$. Reduction modulo such a $p$ and
cancellation of $M/B$ yield
$f_p(B/p)\equiv A\pmod p$. Here $M/B$ is a unit modulo $p$
because $M$ and $B$ are squarefree and $p\mid B$.
For a prime $p\mid t$ outside $B$, the
right side is zero modulo $p$ and reduction instead gives $f_p\equiv0
\pmod p$.

Consequently the exponents have the form
$f_p=r_p+ph_p$ for $p\mid B$, and $f_p=ph_p$ otherwise, with
$h_p\ge0$. (Absent outside primes have $h_p=0$.) Substitution in the
rate sum gives $\lambda=\alpha_0+\sum_p h_p$.
Applying this to the existing $g$ proves that $k=\sum_p h_p$ is a
nonnegative integer. Listing each prime $p$ exactly $h_p$ times gives
\eqref{eq:allscales}. Conversely every product in that formula has
rate $\alpha_0+k=\lambda$. This proves both directions, including
the empty product when $k=0$.
\end{proof}

\begin{corollary}\label{cor:singleton}
An observed homothetic class is a singleton when $k=0$, with sole scale
$t_0$, and infinite when $k\ge1$.
\end{corollary}
\begin{proof}
For $k=0$,~\eqref{eq:allscales} has one value. For $k\ge1$, fix
$k-1$ parameters at 2 and vary the last through primes. The function
$p\mapsto p^p$ is strictly increasing on primes, so this produces
infinitely many distinct scales and first coordinates.
\end{proof}
The denominator condition alone is not an existence assertion for an
arbitrary rational rate. Here existence is supplied by the original
witness; in general the residual $k$ must also be nonnegative.
Theorem~\ref{thm:residues} classifies scales of an admissible form, not
all admissible primitive forms of A230543. For integer rates it recovers
the historical result quoted in Section~\ref{sec:tools}.

\section{The four infinite forms and their specializations}\label{sec:families}
The four infinite forms observed are
\begin{align*}
 Q_4&=(2,9,42,43),\\
 Q_3&=(99,372,1624,1669),\\
 Q_2&=(5202,18003,46294,49943),\\
 Q_H&=(4802517313,5246382600,19023654540,20309803913).
\end{align*}
Their complete first-coordinate families and scale data are shown in
Table~\ref{tab:families}. Prime parameters in all formulas may coincide.

\begin{table}[ht]
\centering\small
\begin{tabular}{clrrrr}
\toprule
Class & First coordinate $n$ & $\lambda$ & $t_0$ & $k$ & Count\\
\midrule
$F_4$ & $2p^pq^qr^rs^s$ & $4$ & $1$ & $4$ & $17$\\
$F_3$ & $99p^pq^qr^r$ & $3$ & $1$ & $3$ & $12$\\
$F_2$ & $25557426p^pq^q$ & $37/17$ & $4913$ & $2$ & $3$\\
$H$ & $4802517313p^p$ & $1$ & $1$ & $1$ & $2$\\
\bottomrule
\end{tabular}
\caption{The four infinite forms. Counts refer only to the 40-term prefix;
the scale for $F_2$ is $17^3p^pq^q$.}\label{tab:families}
\end{table}

\begin{proposition}\label{prop:families}
Each family in Table~\ref{tab:families} produces witnesses for every
choice of its prime parameters, and is the complete homothetic class of
the corresponding ordered form.
\end{proposition}
\begin{proof}
The exact derivative and square calculations are
\begin{align*}
 \D(2)&=1,& \D(9)&=6,& 2^2+9^2+42^2&=43^2,\\
 \D(99)&=75,& \D(372)&=508,& 99^2+372^2+1624^2&=1669^2,\\
 \D(5202)&=6681,& \D(18003)&=7111,&
 5202^2+18003^2+46294^2&=49943^2.
\end{align*}
Thus $1+4\cdot2=9$, $6+4\cdot9=42$;
$75+3\cdot99=372$, $508+3\cdot372=1624$; and
\[
 6681+(37/17)5202=18003,\qquad
 7111+(37/17)18003=46294.
\]
For $Q_H$ the remaining exact checks are
\begin{align*}
 \D(4802517313)&=443865287,\\
 \D(5246382600)&=13777271940,\\
 4802517313^2+5246382600^2+19023654540^2
   &=20309803913^2.
\end{align*}
For $Q_H$, the rate is $1$, so $\D(m)+m=u$ and $\D(u)+u=v$.
Each displayed quadruple has gcd one. The rates for $Q_4,Q_3,Q_H$
are integers, so their cores are 1 and their free counts are 4, 3 and 1.
For $37/17$, $r_{17}=3$, $t_0=17^3$, $\alpha_0=3/17$, and
$k=2$. Theorems~\ref{thm:shape} and~\ref{thm:residues} prove
necessity and sufficiency of all four parameter formulas. Infinitude
follows from Corollary~\ref{cor:singleton}.
\end{proof}

The integer-rate seed chains underlying narrower constructions are
\begin{center}\small
\begin{tabular}{rrrrr}
\toprule
$\lambda$ & $m$ & $T_\lambda(m)$ & $T_\lambda^2(m)$ & Root\\
\midrule
4 & 2 & 9 & 42 & 43\\
3 & 99 & 372 & 1624 & 1669\\
2 & 25557426 & 88448739 & 227442422 & 245369959\\
1 & 102229704 & 353794956 & 909769688 & 981479836\\
\bottomrule
\end{tabular}
\end{center}
For the last two rows, the first and second derivatives of the seed
chain entries are, respectively,
\[
 \D(25557426)=37333887,\quad \D(88448739)=50544944,
\]
\[
 \D(102229704)=251565252,\quad \D(353794956)=555974732.
\]
They verify the displayed transitions directly, and each row satisfies
the three-leg square identity. Formula~\eqref{eq:iterate} therefore
also supplies direct membership proofs for these seed constructions.

\begin{proposition}\label{prop:extraction}
If $g=\prod p^{e_p}$, the greatest divisor of $g$ having an integer
logarithmic derivative is
\begin{equation}\label{eq:extract}
 E=\prod_p p^{p\lfloor e_p/p\rfloor},\qquad
 \Lrate(E)=\sum_p\lfloor e_p/p\rfloor.
\end{equation}
\end{proposition}
\begin{proof}
For a divisor $d=\prod p^{f_p}$ with integer $\sum f_p/p$, multiply
by the squarefree product of its primes and reduce modulo each $p$.
This forces $p\mid f_p$. Conversely these divisibilities make the sum
an integer. The largest admissible exponent not exceeding $e_p$ is
$p\lfloor e_p/p\rfloor$ independently for each prime. Their product
is an admissible divisor divisible by every other admissible divisor.
\end{proof}

For $a(22)$, the full gcd is $78608=2^4 17^3$, with rate $37/17$.
The maximal extraction is $E=16$, with rate 2. Dividing its witness by
16 yields the rate-2 seed row above, because
$\Lrate(78608/16)=3/17$ and
$(78608/16)Q_2=17^3Q_2$. Dividing by 4 instead yields the rate-1
seed $4\cdot17^3 Q_2$. These two seeds are not themselves derivative
witnesses: scaling a form of rate $\lambda$ by $s$, with
$\Lrate(s)=h$, gives a seed of rate $\lambda-h$.
Thus their rates are $37/17-3/17=2$ and
$37/17-20/17=1$, respectively. This explains
\[
 F_1=\{102229704p^p:p\text{ prime}\}\subsetneq F_2:
\]
fix one of $F_2$'s parameters at 2. Strictness follows from the member
$25557426\cdot3^6$, which is not divisible by 4, whereas every member
of $F_1$ is divisible by 4. Both constructions have primitive form
$Q_2$; they are not two independent forms.

In general, if $\Lrate(s)=h$ and the relevant iterates are integral,
the product rule gives
\[
 T_\lambda(sz)=sT_{\lambda+h}(z),\qquad
 T_\lambda^j(sz)=sT_{\lambda+h}^j(z).
\]
This explains the rate correction when a seed is rescaled. Dividing a
witness by its gcd does not preserve its derivative chain without that
correction.

Fixing three parameters of $F_4$ at 2 gives $128p^p$; fixing two
parameters of $F_3$ at 2 gives $1584p^p$. The former is already proved
in the local comparative manuscript on A249106~\cite{comparison}.
Setting all parameters equal gives the diagonals $2p^{4p}$ and
$99p^{3p}$. These are specializations, not additional forms. All four
full families are \emph{parametric control witnesses}: they establish
positive membership and upper bounds, not consecutiveness or global
coverage of the sequence.

\section{Six singleton forms}\label{sec:singletons}
Label by $S_i$ the singleton form observed at index $i$. Their complete
data are in Table~\ref{tab:singletons}; the ordered coordinates are
essential even when the first is not the smallest leg.

\begin{table}[ht]
\centering\scriptsize
\begin{tabular}{clrr}
\toprule
Form & Primitive $Q=(m,u,v,w)$ & $\lambda$ & $t_0$\\
\midrule
$S_2$ & $(1203,404,408,1333)$ & 0 & 1\\
$S_8$ & $(249753,154756,176892,342953)$ & 0 & 1\\
$S_{11}$ & $(722718,843189,281066,1145551)$ & 0 & 1\\
$S_{28}$ & $(942667,194256,536448,1101877)$ & $2/71$ & $71^2$\\
$S_{36}$ & $(258837,169100,283140,419237)$ & $4/19$ & $19^4$\\
$S_{37}$ & $(43108595319,19046592592,38093185200,60598798985)$ & 0 & 1\\
\bottomrule
\end{tabular}
\caption{All six singleton forms. In every row $k=0$ and the associated
term is $a(i)=t_0m$.}\label{tab:singletons}
\end{table}

\begin{proposition}\label{prop:singletons}
Each form in Table~\ref{tab:singletons} has exactly the indicated scale
over all positive integers.
\end{proposition}
\begin{proof}
The exact derivative and square checks are recorded below:
\begin{center}\small
\begin{tabular}{crrr}
\toprule
Form & $\D(m)$ & $\D(u)$ & $m^2+u^2+v^2=w^2$\\
\midrule
$S_2$ & 404 & 408 & 1776889\\
$S_8$ & 154756 & 176892 & 117616760209\\
$S_{11}$ & 843189 & 281066 & 1312287093601\\
$S_{28}$ & 167702 & 530976 & 1214132923129\\
$S_{36}$ & 114608 & 247540 & 175759662169\\
$S_{37}$ & 19046592592 & 38093185200 & 3672214438424437030225\\
\bottomrule
\end{tabular}
\end{center}
The rate-zero rows have $\D(m)=u$, $\D(u)=v$, so $t_0=1$ and
$k=0$. For $S_{28}$,
$u-\D(m)=26554=2m/71$ and
$v-\D(u)=5472=2u/71$; hence $r_{71}=2$, $\alpha_0=2/71$
and $k=0$. For $S_{36}$ the differences are
$54492=4m/19$ and $35600=4u/19$, giving $r_{19}=4$ and $k=0$.
All six quadruples have gcd one. Theorem~\ref{thm:residues} and
Corollary~\ref{cor:singleton} give uniqueness of the scale.
\end{proof}

The corresponding factorizations include
\[
 a(28)=11\cdot17\cdot71^4,
 \qquad a(36)=3\cdot19^6\cdot239,
\]
with full gcds $5041$ and $130321$, respectively, and
\[
 a(37)=3\cdot23\cdot31\cdot53\cdot73\cdot5209
\]
with gcd one. Singleton here is a global statement for a fixed ordered
form, not just within the scanned interval or numerical cap. It does not
exclude a parametrization allowing the form itself to vary. It does
disprove the claim that every term belongs to an infinite class obtained
by scaling its own primitive form.

\section{Exact enumeration and its correctness}\label{sec:algorithm}
The executed program uses C17, OpenMP, GNU unsigned 128-bit intermediates
and GMP integers. Its certified numerical input envelope is
$2\le n\le N=10^{16}$; the completed campaign reaches only
$130000000000$. The input cap is not a claim of practical complete
searchability.

\subsection{Predicate evaluation without family restrictions}
In a block $[\ell,h]$, every integer has a residual initially equal to
$n$ and a derivative accumulator initially zero. An ordinary prime sieve
provides every prime up to $\lfloor\sqrt h\rfloor$. For each such prime,
the segmented loop visits all its multiples in the block, removes its
full exponent $e$ from the residual, and adds $e(n/p)$ to the accumulator.
Afterwards a residual greater than one is prime: if it were composite,
it would have a prime divisor at most $\sqrt h$, already removed.
Adding $n$ divided by that residual completes~\eqref{eq:derivative}.
Hence every input receives its exact first derivative.

The only early arithmetic exclusion is $\D(n)=1$, the prime case
proved impossible in Section~\ref{sec:tools}. For a composite $n$,
either a repeated prime contributes at least 4, or two prime factors
contribute a sum at least $2+3$; thus this exclusion cannot remove a
composite solution. No divisibility, primitivity, shape or congruence
filter is used. The program factors $\D(n)$ exactly to compute
$\D^2(n)$ and then forms the sum of three squares in GMP. Its
\texttt{mpz\_perfect\_square\_p} predicate and
\texttt{mpz\_sqrt} extraction use exact integer arithmetic~\cite{gmp}.

Factorization first tests primality using the seven strong-probable-prime
bases
\[
 2,325,9375,28178,450775,9780504,1795265022.
\]
This is the deterministic test below $2^{64}$ stated in
Fori\v{s}ek and Jan\v{c}ina's Theorem 3~\cite{forisek}, which credits
Sinclair's base set and Izykowski's verification. The operands here are
within that range. Bases are reduced modulo the operand. A zero residue
is skipped; every possible exception divides one of the seven fixed bases.
All 55 distinct divisors at least 2 of those bases were checked against
independent trial-division primality: the implementation accepts the nine
primes and rejects the 46 composites. For all other operands no base
vanishes, and reduction preserves the modular-power tests in the cited
theorem. This finite exceptional-case check is separate from negative
coverage of the sequence. Composite splitting tries small primes through 97,
then a Pollard-rho path~\cite{pollard}, with fixed polynomials
$x^2+c$, $1\le c\le32$, and at most 100000 iterations per polynomial.
If that path fails, exact odd trial division from 101 through
$\lfloor\sqrt x\rfloor$ supplies a deterministic fallback. The
loop guard $p\le x/p$ avoids overflow. Proper factors reduce the
recursive operands; prime leaves are checked, sorted and multiplied back
to the original input. Reconstruction failure is an error, not evidence
of nonmembership. Rho is only a factor-finding optimization; its
heuristic runtime does not enter the proof of exactness.

\subsection{Numeric safety and completion semantics}
For $x\ge2$, if $\Omega(x)$ counts prime factors with multiplicity,
then $\Omega(x)\le\lfloor\log_2x\rfloor$ and
\begin{equation}\label{eq:bound}
 \D(x)=x\sum e_p/p\le\tfrac{x}{2}\lfloor\log_2x\rfloor.
\end{equation}
The corresponding real-logarithm bound is given in
Ufnarovski and \AA hlander's Theorem 9~\cite{ufnarovski}, attributed
there to Barbeau; the floor version above follows directly from
$\Omega(x)\le\lfloor\log_2x\rfloor$.
At $N=10^{16}$ it gives the conservative bounds
\[
 U_1=265000000000000000,
 \qquad U_2=7552500000000000000
\]
for $\D(n)$ and $\D^2(n)$, since
$\lfloor\log_2N\rfloor=53$ and
$\lfloor\log_2 U_1\rfloor=57$. Thus $n$, $\D(n)$ and
$\D^2(n)$ require at most 54, 58 and 63 bits. Factorization operands
are $n$ or $\D(n)$, so 64 prime-factor slots suffice. Modular
multiplication uses a 128-bit intermediate; even the full 64-bit
operand product fits, and the actual envelope is smaller.
The square sum is bounded by $N^2+U_1^2+U_2^2$, a 126-bit integer,
and its floor square root is at most $7557154308997534291<2^{63}$.
GMP nevertheless performs the square computation before conversion;
checked conversions enforce the stated bounds. The sieve requires primes
only through $10^8$ at the full input cap.

OpenMP distributes predicate evaluations with dynamic chunks of 64 and
at most eight threads. Each input owns its flag, taking one of three
states: negative, positive or incomplete. The loop barrier precedes
serial completeness inspection and ordered positive replay. A failed
factorization, out-of-range intermediate or stop request makes the block
incomplete. In particular, a zero second derivative outside the already
excluded prime case is not silently recorded as negative: it causes an
incomplete result. For a composite $n$, $\D(n)\ge4$, so an exact second
derivative is indeed positive.

The durable frontier advances only after a complete block and successful
positive replay. A segment is sealed only when its frontier equals its
upper endpoint. Unsealed partial state is not full-interval evidence;
interrupted blocks are repeated, not treated as intra-input checkpoints.
State writes use flush, file synchronization, atomic rename and directory
synchronization; cooperating commands use a run lock. These mechanisms
preserve completed interval records. Their integrity checks cannot by
themselves establish correctness of each negative predicate evaluation.

\begin{proposition}\label{prop:scan}
In uninterrupted successful execution within the checked numerical
domain, the algorithm considers every integer in its configured interval
and records exactly the witnesses satisfying~\eqref{eq:predicate}.
An interval is admitted as complete only through sealed consecutive
segments reaching its configured endpoint.
\end{proposition}
\begin{proof}
Consecutive blocks partition the interval; the segmented sieve visits all
inputs and computes their first derivatives as proved above. The only
early exclusion is impossible for a witness. Deterministic primality,
exact splitting with fallback and product reconstruction yield the second
derivative, and GMP decides the square condition exactly. Any failed
obligation instead prevents block completion and sealing. Serial replay
records positives in increasing order. Consecutive segment bounds and a
sealed endpoint therefore establish coverage, conditional on correct
execution of these implemented operations.
\end{proof}
This is an algorithmic correctness argument supported by finite tests;
it is not a machine-checked proof of the compiled implementation or a
second negative enumeration.

\section{The 40-term prefix and its validation}\label{sec:extension}
\begin{theorem}[Computational extension, with stated dependencies]
\label{thm:extension}
Assuming the accepted 22-term baseline and correctness of the completed C
enumeration, the prefix in Table~\ref{tab:prefix} is the exact indexed
prefix through $130000000000$. It has 40 terms, with
$a(40)=129667967451$. Its forms comprise 34 members of four infinite
classes and six singleton classes. Positive membership has been checked
independently; full negative coverage has not been independently
enumerated a second time.
\end{theorem}
\begin{proof}[Evidence and structural proof]
The six inclusive run intervals in Table~\ref{tab:coverage} are contiguous
from $408918817$ through $130000000000$. The input count is
$129591081184$, comprising 18 positives and $129591081166$ negatives
across 22 sealed segments.
The sealed parser verifies bounds, configurations, frontier completion,
ordered aggregate equality and joins. Under the stated execution
assumption, Proposition~\ref{prop:scan} excludes omitted terms in that
interval. Appending the sorted positives to the accepted baseline proves
the indicated indices. There is no new local negative enumeration below
the baseline endpoint.

For every row, independent trial division recomputes both derivatives
and an integer square root checks~\eqref{eq:predicate}. Normalization,
rate reconstruction and prime-parameter reconstruction give
Table~\ref{tab:prefix}. The four infinite forms are covered by
Proposition~\ref{prop:families}, and the six remaining forms by
Proposition~\ref{prop:singletons}. Their respective counts sum to
$17+12+3+2=34$ and six. The complete witnesses appear in
Table~\ref{tab:witnesses}.
\end{proof}

\begin{table}[ht]
\centering\small
\begin{tabular}{rr rr}
\toprule
Lower endpoint & Upper endpoint & Segments & Positives\\
\midrule
408918817 & 711541152 & 1 & 1\\
711541153 & 10000000000 & 3 & 7\\
10000000001 & 25728640000 & 2 & 3\\
25728640001 & 47119590400 & 3 & 4\\
47119590401 & 100000000000 & 8 & 2\\
100000000001 & 130000000000 & 5 & 1\\
\bottomrule
\end{tabular}
\caption{Completed inclusive coverage after the accepted baseline.}
\label{tab:coverage}
\end{table}

Table~\ref{tab:coverage} is the primary coverage record; the following
bounded cross-checks do not replace its full negative enumeration.
The independent campaign windows have a union of 10763 distinct inputs,
with 11 positives and 10752 negatives. The final window is
$129667966939$--$129667967963$, containing 1025 inputs, one positive and
1024 negatives. Earlier window executions are inherited and reconciled
with their original reports, not represented as fresh executions. All
40 positives were rechecked during manuscript preparation. The separate
code-diagnostic union covers 33827 distinct inputs with 30 positives and
33797 negatives; it is not added to the campaign-window count.

These checks have different meanings. The independent Python oracle uses
trial division and integer roots, unlike the C sieve, primality and rho
route. Native replay checks state and recomputes positives, not all
negatives. Archive integrity checks bind recorded files but are not
mathematical certificates. The final auditor was run with ordinary
Python, with \texttt{python3 -O}, and against the live retained state.
Those reports agree, but they share the same audit implementation. Retained
finite code tests include forced fallback and incomplete-block behavior,
state corruption/recovery and parallel configurations; they support the
implementation argument without certifying every input independently.

Initial family controls were known before the runs that assigned their
indices. In contrast, the form $Q_H$ first appeared in the run reaching
$a(33)=19210069252$, with parameter 2. Its parameter-3 witness was then
known before the final run indexed it at 40. Thus the last witness
\[
 (129667967451,141652330200,513638672580,548364705651)
\]
confirms a previously observed form rather than introducing an eleventh.

\subsection{Recorded resource use}
The recorded environment is an AMD Ryzen 9 7940HS system running Linux
Mint 22.3 (kernel 6.17.0-42) on x86\_64 with glibc 2.39. GCC 13.3.0 compiled the program
with \texttt{-O3}, \texttt{-std=c17}, \texttt{-fopenmp}, strict warning
flags and \texttt{-Werror}; it links GMP 6.3.0, OpenSSL 3.0.13 and
GNU OpenMP. All campaign configurations use eight threads.
The sum of recorded continuation-call wall times is about 62022 seconds,
or 17 hours 14 minutes. The final $100$--$130$ billion continuation uses
about 15946 seconds, or 4 hours 26 minutes, in five calls of approximately
52--54 minutes each; its maximum recorded resident set is 23536 KiB.
These times exclude development, reviews and audits and are not calendar
elapsed time. No full-cap timing or unmeasured parallel speedup is inferred.

The retained code, witnesses, raw interval records and independent
validators are the basis of the local result. The published repository
and release series in~\cite{project} supply those artifacts. Actual initial
release payloads from GitHub and Zenodo were downloaded and validated with
the bounded public suite; this does not repeat the complete negative scan.

\section{Finite auxiliary studies and related predicates}\label{sec:finite}
\subsection{A bounded seed rectangle}
For $1\le\lambda\le8$ and $1\le m\le20000$, the seed test checks
positivity of the iterates and whether
$m^2+T_\lambda(m)^2+T_\lambda^2(m)^2$ is an integer square.
The retained study covers exactly 160000 pairs, with the following
15 positives:
\begin{center}\small
\begin{tabular}{cl}
\toprule
$\lambda$ & Positive $m$\\
\midrule
1 & 128, 864, 1584, 5832, 10692\\
2 & 32, 216, 396, 1458, 2673\\
3 & 8, 54, 99, 6250\\
4 & 2\\
5--8 & none in this rectangle\\
\bottomrule
\end{tabular}
\end{center}
These seeds normalize to just $Q_4$ and $Q_3$. Trial division checked
the rectangle, and SymPy independently rechecked its positives; the
159985 negatives were not enumerated by a second independent
implementation. Neither completeness outside the rectangle nor
completeness among rational-rate primitive forms follows.

\subsection{The prime pattern behind index 28}
Since $187=11\cdot17$ and $\D(187)=28$, the product rule, including
shared-prime cases $p=11,17$, gives for every prime $p$
\begin{equation}\label{eq:pattern}
 C_p=28p+748,\quad \D(187p^4)=p^3C_p,\quad
 \D^2(187p^4)=p^2\bigl(p\D(C_p)+3C_p\bigr).
\end{equation}
The second equality follows by differentiating $p^3C_p$, without any
coprimality condition. The retained records cover all 168 primes
$p\le1000$ and have exactly one positive, $p=71$, giving $a(28)$.
Their read-only replay validates the stated finite screen. No assertion
of uniqueness over all primes follows from it.

\subsection{Actual mathematical correlations}
A003415 supplies $\D(n)$ and A068346 supplies
$\D(\D(n))$~\cite{derivatives}. A051674 supplies exactly the
rate-one scale factors, so products of its values, together with a
fixed rational-rate core, implement~\eqref{eq:allscales}. This is a
structural relation, not equality of sequence indices.

A210503 requires $n,\D(n)$ to be legs of a primitive Pythagorean
triple. A211176 instead requires $n$ to be the hypotenuse of a
Pythagorean triple with $\D(n)$ a leg, without a primitivity
requirement~\cite{triples}. Deleting $\D^2(n)$ from
\eqref{eq:predicate} proves neither condition. The primitive-quadruple
catalogues A096907--A096909 and A097263--A097265 project sorted legs
or parity-selected legs, ordered by diagonal and then by the other
coordinates~\cite{catalogues}. Our $Q$ must retain derivative order;
sorting it forgets which coordinate is the original integer. A097266
counts primitive quadruples by odd diagonal, and A360946 counts
quadruples by inradius~\cite{counts}. Such catalogues can identify a
geometric form, but do not furnish a correspondence with the increasing
$n$ indices here or settle an unsearched interval.

\section{The first unresolved index and open questions}\label{sec:endpoint}
By $F_4$ with parameters $2,5,5,5$, there is an exact witness
\begin{align*}
 n&=244140625000=2\cdot2^2\cdot5^{15},\\
 \D(n)&=1098632812500,\\
 \D^2(n)&=5126953125000,\\
 R&=5249023437500.
\end{align*}
Together with completed coverage this gives, under the indexing
dependencies of Theorem~\ref{thm:extension},
\begin{equation}\label{eq:endpoint}
 130000000000<a(41)\le244140625000.
\end{equation}
The upper witness is not an established value of $a(41)$ and is not in
the indexed 40-term table. Negative coverage of
$130000000001$--$244140624999$, or another exact ordering proof,
is still required to identify that index.

A finite parameter comparison finds this to be the smallest member above
the completed endpoint among the four observed infinite classes. To make
that comparison exhaustive within those classes up to the bound, it
suffices to use primes $2,3,5,7$: already $11^{11}$ exceeds the bound
even before multiplication by a seed or other parameters. This does not
exclude a different primitive form or an intermediate singleton.

The sequence is infinite by $F_4$, but infinitude of the set of
\emph{distinct primitive forms} remains open in this study. Are there
infinitely many admissible singleton forms? How can admissible primitive
forms be characterized before their first witness is found? The ten
observed forms and two bounded screens do not justify conjectural
answers to these questions. A universal cover by variable-form families
would first require a precise definition of the permitted
parametrizations. Saying that every term belongs to some homothetic
class is merely the definition, while saying that every such class is
infinite is already false.

\section{Conclusions}
The local computation adds 18 indexed terms to the dated accepted baseline,
yielding 40 terms through $130000000000$, with the negative-enumeration dependency
and independent validation scopes stated explicitly. The fixed-form
criterion separates complete scale classification from the unresolved
classification of all primitive forms. It proves four infinite classes,
accounts for their specializations, and proves six singleton classes
over the full positive-integer scale domain. The next index has an exact
upper witness, but remains unresolved. Finite seed and pattern studies
give additional evidence without extending their bounds into global
claims.

\appendix
\section{Complete indexed and witness tables}\label{sec:tables}
In Table~\ref{tab:prefix}, $g$ is the full quadruple gcd and the prime
tuple gives the free parameters of the corresponding complete family.
For $S_i$ the tuple is empty and the sole scale is the indicated $g$.
The $F_2$ scale includes its fixed core $17^3$. The horizontal break
after 22 distinguishes the accepted baseline from the local extension.
The witness table supplies every derivative and root; both tables are
part of this monolithic manuscript.

\begingroup\small
\begin{longtable}{rrcrl}
\caption{All 40 indexed terms, forms, full gcds and free prime parameters.}\label{tab:prefix}\\
\toprule
$i$ & $a(i)$ & Class & $g$ & Parameters\\
\midrule
\endfirsthead
\multicolumn{5}{c}{Table \thetable\ (continued)}\\
\toprule
$i$ & $a(i)$ & Class & $g$ & Parameters\\
\midrule
\endhead
\bottomrule
\endfoot
1 & 512 & $F_4$ & 256 & $(2,2,2,2)$\\
2 & 1203 & $S_{2}$ & 1 & $\varnothing$\\
3 & 3456 & $F_4$ & 1728 & $(2,2,2,3)$\\
4 & 6336 & $F_3$ & 64 & $(2,2,2)$\\
5 & 23328 & $F_4$ & 11664 & $(2,2,3,3)$\\
6 & 42768 & $F_3$ & 432 & $(2,2,3)$\\
7 & 157464 & $F_4$ & 78732 & $(2,3,3,3)$\\
8 & 249753 & $S_{8}$ & 1 & $\varnothing$\\
9 & 288684 & $F_3$ & 2916 & $(2,3,3)$\\
10 & 400000 & $F_4$ & 200000 & $(2,2,2,5)$\\
11 & 722718 & $S_{11}$ & 1 & $\varnothing$\\
12 & 1062882 & $F_4$ & 531441 & $(3,3,3,3)$\\
13 & 1948617 & $F_3$ & 19683 & $(3,3,3)$\\
14 & 2700000 & $F_4$ & 1350000 & $(2,2,3,5)$\\
15 & 4950000 & $F_3$ & 50000 & $(2,2,5)$\\
16 & 18225000 & $F_4$ & 9112500 & $(2,3,3,5)$\\
17 & 33412500 & $F_3$ & 337500 & $(2,3,5)$\\
18 & 105413504 & $F_4$ & 52706752 & $(2,2,2,7)$\\
19 & 123018750 & $F_4$ & 61509375 & $(3,3,3,5)$\\
20 & 225534375 & $F_3$ & 2278125 & $(3,3,5)$\\
21 & 312500000 & $F_4$ & 156250000 & $(2,2,5,5)$\\
22 & 408918816 & $F_2$ & 78608 & $(2,2)$\\
\midrule
23 & 711541152 & $F_4$ & 355770576 & $(2,2,3,7)$\\
24 & 1304492112 & $F_3$ & 13176688 & $(2,2,7)$\\
25 & 2109375000 & $F_4$ & 1054687500 & $(2,3,5,5)$\\
26 & 2760202008 & $F_2$ & 530604 & $(2,3)$\\
27 & 3867187500 & $F_3$ & 39062500 & $(2,5,5)$\\
28 & 4751984347 & $S_{28}$ & 5041 & $\varnothing$\\
29 & 4802902776 & $F_4$ & 2401451388 & $(2,3,3,7)$\\
30 & 8805321756 & $F_3$ & 88942644 & $(2,3,7)$\\
31 & 14238281250 & $F_4$ & 7119140625 & $(3,3,5,5)$\\
32 & 18631363554 & $F_2$ & 3581577 & $(3,3)$\\
33 & 19210069252 & $H$ & 4 & $(2)$\\
34 & 26103515625 & $F_3$ & 263671875 & $(3,5,5)$\\
35 & 32419593738 & $F_4$ & 16209796869 & $(3,3,3,7)$\\
36 & 33731896677 & $S_{36}$ & 130321 & $\varnothing$\\
37 & 43108595319 & $S_{37}$ & 1 & $\varnothing$\\
38 & 59435921853 & $F_3$ & 600362847 & $(3,3,7)$\\
39 & 82354300000 & $F_4$ & 41177150000 & $(2,2,5,7)$\\
40 & 129667967451 & $H$ & 27 & $(3)$\\
\end{longtable}
\endgroup

\begingroup\small
\begin{longtable}{rrrrr}
\caption{Complete positive witnesses for the indexed prefix.}\label{tab:witnesses}\\
\toprule
$i$ & $n$ & $\D(n)$ & $\D^2(n)$ & $R$\\
\midrule
\endfirsthead
\multicolumn{5}{c}{Table \thetable\ (continued)}\\
\toprule
$i$ & $n$ & $\D(n)$ & $\D^2(n)$ & $R$\\
\midrule
\endhead
\bottomrule
\endfoot
1 & 512 & 2304 & 10752 & 11008\\
2 & 1203 & 404 & 408 & 1333\\
3 & 3456 & 15552 & 72576 & 74304\\
4 & 6336 & 23808 & 103936 & 106816\\
5 & 23328 & 104976 & 489888 & 501552\\
6 & 42768 & 160704 & 701568 & 721008\\
7 & 157464 & 708588 & 3306744 & 3385476\\
8 & 249753 & 154756 & 176892 & 342953\\
9 & 288684 & 1084752 & 4735584 & 4866804\\
10 & 400000 & 1800000 & 8400000 & 8600000\\
11 & 722718 & 843189 & 281066 & 1145551\\
12 & 1062882 & 4782969 & 22320522 & 22851963\\
13 & 1948617 & 7322076 & 31965192 & 32850927\\
14 & 2700000 & 12150000 & 56700000 & 58050000\\
15 & 4950000 & 18600000 & 81200000 & 83450000\\
16 & 18225000 & 82012500 & 382725000 & 391837500\\
17 & 33412500 & 125550000 & 548100000 & 563287500\\
18 & 105413504 & 474360768 & 2213683584 & 2266390336\\
19 & 123018750 & 553584375 & 2583393750 & 2644903125\\
20 & 225534375 & 847462500 & 3699675000 & 3802190625\\
21 & 312500000 & 1406250000 & 6562500000 & 6718750000\\
22 & 408918816 & 1415179824 & 3639078752 & 3925919344\\
\midrule
23 & 711541152 & 3201935184 & 14942364192 & 15298134768\\
24 & 1304492112 & 4901727936 & 21398941312 & 21991892272\\
25 & 2109375000 & 9492187500 & 44296875000 & 45351562500\\
26 & 2760202008 & 9552463812 & 24563781576 & 26499955572\\
27 & 3867187500 & 14531250000 & 63437500000 & 65195312500\\
28 & 4751984347 & 979244496 & 2704234368 & 5554561957\\
29 & 4802902776 & 21613062492 & 100860958296 & 103262409684\\
30 & 8805321756 & 33086663568 & 144442853856 & 148445272836\\
31 & 14238281250 & 64072265625 & 299003906250 & 306123046875\\
32 & 18631363554 & 64479130731 & 165805525638 & 178874700111\\
33 & 19210069252 & 20985530400 & 76094618160 & 81239215652\\
34 & 26103515625 & 98085937500 & 428203125000 & 440068359375\\
35 & 32419593738 & 145888171821 & 680811468498 & 697021265367\\
36 & 33731896677 & 22037281100 & 36899087940 & 54635385077\\
37 & 43108595319 & 19046592592 & 38093185200 & 60598798985\\
38 & 59435921853 & 223334979084 & 974989263528 & 1002005591643\\
39 & 82354300000 & 370594350000 & 1729440300000 & 1770617450000\\
40 & 129667967451 & 141652330200 & 513638672580 & 548364705651\\
\end{longtable}
\endgroup

\section*{Data availability}
The indexed prefix, documentary witness table, executed source and build
instructions, sealed coverage and timing records, finite-study data and
package-local validation scripts are publicly available in the GitHub
repository and Zenodo release series cited in~\cite{project}. The archival identifier cited there is the concept DOI for the release
series; individual deposited releases have separate version-specific DOIs.
The package supplies this monolithic source and its matching PDF.
Reproduction of the coverage audit requires the supplied raw records and
documented validator dependencies in addition to the positive witness data;
the manuscript alone does not replace those artifacts. Historical development
logs are archival evidence, not newly repeated sanitizer certification.
Private reviewer conversations and third-party full texts are excluded.

\section*{AI use disclosure}
OpenAI Codex and ChatGPT, Anthropic Claude, Google Gemini, xAI Grok,
Alibaba Qwen and Meta AI contributed in the roles described below.
The retained code reports identify ChatGPT as GPT-5.6 Sol, Claude as
Opus 5.5, Gemini as 3.0 Pro, Grok as 4.7, Qwen as 3.8 and Meta AI as
Muse Spark 1.1. These are reported provenance, not independently
authenticated identities. Other contributions lack sufficient provider or
model metadata for a reliable family attribution. Same-provider
contributions are not assumed independent.

\noindent\textbf{Mathematical analysis.}
OpenAI Codex assisted with source reading, exact-arithmetic checks,
family analysis and formulation of the proofs. Historical sources and
local evidence were checked separately from generated explanations.

\noindent\textbf{Algorithm and code.}
OpenAI Codex assisted with algorithm design and implementation. The
retained code-review process includes contributions from all the families
named above. Frozen review and local validation records document their
scope separately from the executed campaign.
Reviewer execution claims were not substituted for local validation.

\noindent\textbf{Manuscript.}
OpenAI Codex, based on GPT-6, assisted with drafting and internal
six-pass checking and revision of this manuscript. Nine first-round
manuscript reports were subsequently supplied by the author. Reported
review systems include the six providers named above and Z.ai GLM;
two reports do not establish a model/provider family. One report
explicitly lacked attachment access. The reports were adjudicated against
the sources and retained evidence; their execution claims were not
substituted for local validation. The historical code review remains
separate from manuscript review. Numerical claims depend on the retained computational
artifacts and local arithmetic checks, not on conversational assertions.
Carlo Corti assumes full scientific responsibility for the mathematical
statements, algorithm, implementation, numerical results, validation boundaries
and conclusions.

\begingroup
\raggedright
\clearpage
\begin{thebibliography}{99}
\bibitem{sequence}
P. P. Lava, OEIS A230543 (2013); description: numbers $n$ such that
$n$, $n'$ and $n''$ are the legs of a Pythagorean quadruple, with terms from G. Resta
(2013) and R. Chandler (2016--2017).
\url{https://oeis.org/A230543}.
Baseline acquired October 6, 2026; current endpoint not reverified.

\bibitem{ufnarovski}
V. Ufnarovski and B. \AA hlander, How to differentiate a number,
\emph{Journal of Integer Sequences} \textbf{6} (2003), Article 03.3.4.
\url{https://cs.uwaterloo.ca/journals/JIS/VOL6/Ufnarovski/ufnarovski.html}.

\bibitem{kovic}
J. Kovi\v{c}, The arithmetic derivative and antiderivative,
\emph{Journal of Integer Sequences} \textbf{15} (2012), Article 12.3.8.
\url{https://cs.uwaterloo.ca/journals/JIS/VOL15/Kovic/kovic4.pdf}.

\bibitem{book}
G. Balzarotti and P. P. Lava, \emph{La derivata aritmetica}, Hoepli,
2013. ISBN 9788820358648.
\url{https://www.hoepli.it/libro/la-derivata-aritmetica-/9788820358648.html}.
Publisher metadata consulted; full text not consulted.

\bibitem{sandit}
G. Balzarotti and P. P. Lava, \emph{La derivata aritmetica}, Sandit Libri.
ISBN 9788869287480. Publisher metadata consulted; publication year
not independently verified and full text not consulted.

\bibitem{overview}
T. Tossavainen, P. Haukkanen, J. K. Merikoski and M. Mattila,
We can differentiate numbers, too,
\emph{The College Mathematics Journal} \textbf{55}, no. 2 (2024), 100--108.
DOI: 10.1080/07468342.2023.2268494.
\url{https://www.researchgate.net/publication/375299116_We_Can_Differentiate_Numbers_Too}.
Full author-uploaded text consulted October 9, 2026.

\bibitem{derivatives}
N. J. A. Sloane and R. K. Guy, OEIS A003415, arithmetic derivative;
R. Zumkeller, OEIS A068346, second arithmetic derivative;
A. Auel, OEIS A051674, $p^p$ for prime $p$, with T. D. Noe's
fixed-point interpretation (2004).
\url{https://oeis.org/A003415}; \url{https://oeis.org/A068346};
\url{https://oeis.org/A051674}. Definitions consulted October 9, 2026.

\bibitem{comparison}
C. Corti, \emph{An Exact Extension and an Infinite Family for OEIS A249106},
local manuscript, version 4 (2026), section on A230543.
Public availability not verified.

\bibitem{forisek}
M. Fori\v{s}ek and J. Jan\v{c}ina, Fast primality testing for integers that
fit into a machine word, in \emph{Student Research Forum Papers and
Posters at SOFSEM 2015}, CEUR Workshop Proceedings \textbf{1326}
(2015), 20--30.
\url{https://ceur-ws.org/Vol-1326/020-Forisek.pdf}.

\bibitem{pollard}
J. M. Pollard, A Monte Carlo method for factorization,
\emph{BIT} \textbf{15} (1975), 331--334.
DOI: 10.1007/BF01933667.
\url{https://pages.cs.wisc.edu/~cs812-1/pollardrho.pdf}.

\bibitem{gmp}
GNU MP project, GNU MP 6.3.0 manual, Section 5.8, Root extraction functions.
\url{https://gmplib.org/manual/Integer-Roots}. Consulted October 9, 2026.

\bibitem{triples}
P. P. Lava, OEIS A210503, arithmetic-derivative primitive Pythagorean
triples, and OEIS A211176, arithmetic-derivative hypotenuse condition
(2013). \url{https://oeis.org/A210503}; \url{https://oeis.org/A211176}.
Definitions consulted October 9, 2026.

\bibitem{catalogues}
R. Chandler, OEIS primitive Pythagorean quadruple projections
A096907--A096909 and A097263--A097265 (2004).
\url{https://oeis.org/A096907}; \url{https://oeis.org/A096908};
\url{https://oeis.org/A096909}; \url{https://oeis.org/A097263};
\url{https://oeis.org/A097264}; \url{https://oeis.org/A097265}.
Definitions consulted October 9, 2026.

\bibitem{counts}
R. Chandler, OEIS A097266, Number of primitive Pythagorean quadruples
with diagonal $2n+1$;
M.-\'{A}. P\'{e}rez Garc\'{i}a-Ortega, OEIS A360946,
Number of Pythagorean quadruples with inradius $n$ (2023).
\url{https://oeis.org/A097266}; \url{https://oeis.org/A360946}.
Definitions consulted October 9, 2026.

\bibitem{project}
C. Corti, \emph{An Exact Local Extension and Homothetic Classes for
OEIS A230543}, GitHub repository and Zenodo archive, 2026.
\url{https://github.com/carcorti/A230543};
\url{https://doi.org/10.5281/zenodo.23311392}.
The DOI identifies the archived release series; each deposited release has
a separate version-specific DOI.
\end{thebibliography}
\endgroup
\end{document}

